%%=============================================================================
%% Conclusie
%%=============================================================================

\chapter{Conclusie}%
\label{ch:conclusie}

\section{Conclusie}
\label{sec:conclusie}

Deze bachelorproef onderzocht in welke mate zes large language models betrouwbare en technisch correcte Python code genereren voor 35 representatieve data analyse en ETL taken (630 individuele evaluaties). De evaluatie gebeurde volledig geautomatiseerd op acht dimensies: syntaxis, style, beveiliging, uitvoerbaarheid, performantie, schaalbaarheid, onderhoudbaarheid (Radon) en functionele correctheid.

De centrale onderzoeksvraag kan als volgt worden beantwoord: de technische correctheid en betrouwbaarheid van LLM-gegenereerde Python code verschilt aanzienlijk tussen modellen, maar geen enkel model behaalde consistent sterke resultaten over alle kwaliteitscriteria. De algemene rangschikking op basis van de gemiddelde kwaliteitsscore over de zes codekwaliteitsdimensies (syntaxis, style, beveiliging, uitvoerbaarheid, performantie en onderhoudbaarheid, dus exclusief functionele correctheid) is: \texttt{deepseek-coder-v2} (90,3\%), \texttt{qwen2.5-coder} (90,2\%), \texttt{gemma4} (86,5\%), \texttt{phi4-mini} (85,4\%), \texttt{llama3} (81,7\%) en \texttt{mistral} (79,6\%).

Een opvallende bevinding is het significante verschil tussen codekwaliteit en functionele correctheid. Terwijl alle modellen uitstekende scores behaalden op syntaxis (95-100\%), beveiliging (100\%) en onderhoudbaarheid (>95\%), bleef de functionele correctheid tussen 29,5\% en 58,1\%. Dit wijst erop dat LLM's weliswaar syntactisch correcte code kunnen genereren, maar dat het implementeren van de correcte ETL-logica aanzienlijk moeilijker blijkt.

De style score vertoonde de grootste spreiding (54,3\% tot 100\%). Opvallend is dat de lage score van Gemma 4 (54,3\%) niet wijst op fundamenteel slechte code, maar op triviale schendingen van de PEP-8 stijlrichtlijnen: ontbrekende nieuwe lijn op het einde van het bestand, te lange regels, witruimte op blanco lijnen en ongebruikte imports. Deze problemen kunnen door Ruff automatisch worden gecorrigeerd en beïnvloeden de correctheid van de code niet. Qwen 2.5 Coder was het enige model dat op alle 105 samples perfecte PEP-8 compliance behaalde (100\%).

Uit de opsplitsing per moeilijkheidsniveau (zie hoofdstuk Resultaten) blijkt dat de functionele correctheid voor eenvoudige taken varieerde van 44,4\% tot 75,0\%, maar voor gemiddelde taken daalde naar 20,0\% tot 55,0\%. Alle zes modellen presteerden lager op gemiddelde dan op eenvoudige taken, wat aantoont dat taken waarbij meerdere bewerkingen moeten worden gecombineerd — zoals inlezen, transformeren, valideren en wegschrijven — systematisch vaker tot fouten leiden. De beveiligingsanalyse leverde een 100\% score op voor alle modellen, wat te verklaren is doordat de ETL-taken in deze benchmark doorgaans geen beveiligingsgevoelige operaties bevatten zoals netwerkverkeer, authenticatie of cryptografie.

Ter vergelijking: \textcite{Bayram2025} rapporteren op HumanEval, een benchmark van eenvoudige ge\"isoleerde Python functies van gemiddeld 7,2 regels code, functionele scores tussen 72\% en 95\% voor grote commerci\"ele modellen zoals GPT-4 Omni en Claude Sonnet 4 (beide met 100+ miljard parameters). In dit onderzoek worden open source modellen van 3,8 tot 16 miljard parameters gebruikt, die uitvoerbaar zijn op gewone consumentenhardware. De complexiteit van de opdrachten verschilt bovendien aanzienlijk: een ETL-taak met inlezen, transformeren, valideren en wegschrijven over meerdere stappen en datasets is fundamenteel moeilijker dan het implementeren van een ge\"isoleerde stringbewerking of wiskundige functie. Een functionele score van 57,1\% voor de best presterende modellen is daarom realistisch en vergelijkbaar met wat van kleinere open source modellen op deze complexe taken mag worden verwacht.

Ook op vlak van schaalbaarheid (zie Tabel~\ref{tab:pass_rates} en Figuur~\ref{fig:scalability}) waren er duidelijke verschillen tussen de modellen. Qwen 2.5 Coder behaalde de hoogste schaalbaarheidsscore (80,6\%), gevolgd door Gemma 4 (79,9\%) en DeepSeek Coder V2 (77,2\%). Mistral scoorde het zwakst (52,7\%). Uit de analyse blijkt dat alle modellen voor bepaalde taken een score van 0\% noteerden, wat betekent dat de code wel correct werkte op kleine datasets maar faalde bij opschaling naar grotere invoer. Dit wijst erop dat geen enkel model standaard schaalbare code genereert; expliciete instructies voor geheugenefficiëntie of chunking blijven noodzakelijk.

De meerwaarde van dit onderzoek ligt in een concreet, volledig geautomatiseerd evaluatiekader voor studenten Toegepaste Informatica aan HOGENT. De bevindingen tonen aan dat generatieve AI nuttig kan zijn als ondersteunende programmeerhulp, maar niet zonder controle mag worden gebruikt. Zelfs de best presterende modellen halen geen volledig overtuigend niveau, met name door de systematische zwakte op functionele correctheid. De resultaten bevestigen de nood aan een kritische, testgebaseerde aanpak bij het gebruik van generatieve AI in programmeeronderwijs.

Tegelijk roept het onderzoek nieuwe vragen op, bijvoorbeeld over de impact van prompt engineering, andere temperatuurinstellingen of een grotere variëteit aan ETL-scenario's. Dit biedt interessante mogelijkheden voor vervolgonderzoek.

Samengevat toont deze bachelorproef aan dat LLM's vandaag al nuttig zijn als codegenererende assistent, maar dat hun output voor data analyse en ETL nog onvoldoende betrouwbaar is om zonder controle te gebruiken. Voor studenten ligt de concrete meerwaarde in het bewust leren evalueren van AI-gegenereerde code in plaats van die automatisch te vertrouwen.
